\chapter{文献综述}
\label{ch:literature}

\section{作为可执行层级的行为树}

行为树会重复评估一棵有根层级，并传播 \textsc{success}、\textsc{failure} 或 \textsc{running} 三种状态。Sequence 通常在第一个失败或运行中的子节点停止；Selector 在第一个成功或运行中的子节点停止；叶节点读取状态或执行动作。这个小型状态接口使复杂逻辑能够组合，并能明确表达中断策略\parencite{colledanchise2018bt}。

三状态接口看似简单，却包含关键实现选择。运行中的 Sequence 可以记住当前索引，而 Reactive Selector 需要重新检查高优先级条件并抢占较低优先级分支；Random Selector 在子节点运行期间应保持同一次选择；Parallel 需要成功和失败阈值；Repeat 与 Wait 需要在重启或中断时清除记忆。因此，编辑器显示和运行语义不能完全分离：子节点顺序、Decorator 归属和节点身份都会影响执行，而且必须正确序列化。

黑板在感知、条件和动作之间提供共享键值状态。Unreal Engine 的工作流说明了把黑板、Decorator 和运行观察结合的价值\parencite{unrealbt}。但本项目没有必要重现商业系统的全部能力。合理基线是：系统能够表达有意义的决策、调用项目 Actor 方法、验证黑板键，并解释分支为何运行或被拒绝。

\section{节点--连线树与扩展性}

节点--连线图适合行为树，因为边能直接表示祖先关系和同级分组。经典整洁树算法追求层级一致、子树不重叠且相对紧凑的布局\parencite{reingold1981tidier}。SpaceTree 又把节点--连线结构与动态布局、缩放和选择性展开结合，其评估说明表示和交互会共同影响拓扑理解与重访任务\parencite{plaisant2002spacetree}。

问题在于规模。卡片占据面积，随着宽度和深度增加，显式连线也会变密。Song 等比较传统、列表和多列视图，发现多列界面在其大扇出层级的浏览和理解任务中更快，总体偏好也更高\parencite{song2010largefanout}。这一结果适用于包含大量备选行为的 Selector 或 Sequence，但不能直接照搬：论文中的子节点按字母排序，而行为树的横向顺序可能决定优先级。因此多列行为树必须明确显示顺序，且不能在布局时静默修改资源。

Bae 与 Watson 为大型分层图提出 Quilts 表示\parencite{bae2011quilts}。其价值不一定是要求行为树完全改用矩阵，而是证明路径任务可能更适合另一种补充表示。本项目因此保留可编辑节点图，同时加入紧凑文字路径和强活动路径编码，以针对运行时路径定位任务。

\section{焦点、上下文与概览}

Furnas 用兴趣度函数形式化了广义鱼眼视图：综合元素先验重要性和与焦点的距离，决定显示显著程度\parencite{furnas1986fisheye}。Lamping 等使用双曲几何在有限区域显示大型层级，对焦点进行放大并压缩远处上下文\parencite{lamping1995hyperbolic}。行为树编辑器可以借此读取鼠标所在节点，同时保留周边结构。

但是焦点+上下文并不天然优于其他方法。畸变会移动目标并增加点击难度，缩放变化也可能破坏心理地图。Cockburn 等区分概览+细节、缩放和焦点+上下文，并指出效果取决于任务、实现和切换成本\parencite{cockburn2008review}。因此本项目把鱼眼最大倍率限制为 1.20，改变卡片内部表示而不是施加不稳定的整体图变换，以节点中心补偿位置，并在关闭时恢复全部几何。同时提供独立固定小地图，使两类技术可以比较或组合。

语义缩放与几何缩放不同。几何缩放改变物理大小，语义缩放改变显示属性。低细节卡片保留类型颜色、图标和短标题；高细节卡片展示方法、参数、描述和 Decorator。它以信息完整性换取密度，却不改变执行图几何。此前插件曾把滚轮与临时节点缩放耦合，造成滚动时节点先缩小再放大；分离信息可见性和图几何后，该问题被消除。

\section{复杂度管理与心理地图}

当不需要同时查看全部节点时，展开--折叠和隐藏--显示可以降低复杂度。Dogrusoz 等提出大型网络的增量复杂度管理方法，包括展开、折叠和布局更新\parencite{dogrusoz2018complexity}。该研究与心理地图原则一致：如果一次编辑引起大量无关节点移动，用户就必须重新学习位置\parencite{misue1995mentalmap}。行为树折叠应保留折叠根，报告隐藏内容，并尽量不移动无关分支。

子树折叠和子树聚焦解决不同问题。折叠用摘要替代后代，同时保留树的其他部分；聚焦只显示目标分支和必要上下文。二者都可能隐藏运行节点，因此插件提供隐藏节点数量、两层摘要、路径指示以及显式展开/聚焦控制，而不是无提示地删除视图信息。

搜索高亮与非活动分支淡化属于较柔和的复杂度管理：几何保持不变，只降低非相关内容显著性。这符合 Shneiderman 的“先概览、再缩放和过滤、按需细节”原则\parencite{shneiderman1996eyes}。形状/图标和色盲友好配色则提供冗余通道，避免只用颜色表达节点类型。

\section{运行时解释与失败定位}

实时调试为静态层级增加时间维度。只高亮叶节点不能解释其到达路径；高亮所有历史节点又会混淆当前状态。本项目记录有序的 Root 到叶节点链条、逐节点状态和当前失败解释。非活动分支可以淡化，独立路径摘要在图缩小时仍可阅读。

Decorator 让失败解释更加重要。Action 本身可能正确，却因为黑板比较、冷却或时间限制而从未执行。在所属节点上显示 Decorator 摘要并附加失败原因，可以把运行证据连接回编写表示；实时黑板进一步显示 Key、运行时类型、值和 Schema 状态。这体现了分层图研究的任务导向原则：调试时用户往往在问“为什么这个分支没有运行”，而不是“完整拓扑是什么”。

\section{研究空白与设计启示}

文献支持多个单项机制，却没有规定如何把它们集成到一个可执行 Godot 行为树编辑器中。本文据此采用四项原则：保留节点--连线拓扑和明确同级顺序作为基线；分离几何放大、语义细节和结构过滤；在折叠与布局中维护空间上下文，并在隐藏细节时提供概览或路径表示；最后，把几何效果和软件正确性与人的任务表现分开评估。最后一项尤其重要：面积或可见节点减少只能证明表示更紧凑，不能证明更容易理解。
